\chapter{Spherical harmonics and higher spectral gaps}\label{sec:spherical-harmonics-local}
The round-sphere spectrum enters both constrained CMC stability and Gaussian stability. Here the harmonic-polynomial decomposition is proved algebraically before passing to the Sobolev energy. The reference is \cite{ABR01}, second edition, author PDF dated July 17, 2020, Chapter 5, especially Proposition 5.5, Theorem 5.7, Proposition 5.8 and Theorem 5.12. Our sphere has dimension $n$, so the ambient dimension in that source is $n+1$. Uniform polynomial approximation on a compact sphere (Stone--Weierstrass), density of continuous functions in $L^2$ for its finite Radon measure, and completeness of $H^1$ remain background inputs.

\begin{definition}[Homogeneous and harmonic polynomial spaces]\label{def:homogeneous-harmonic-spaces}
For $n\geq1$ and $k\geq0$, let $\mathcal P_k$ be the real homogeneous polynomials of degree $k$ on $\mathbb R^{n+1}$, and let $\mathcal H_k=\ker(\Delta:\mathcal P_k\to\mathcal P_{k-2})$. Set $\mathcal P_j=\{0\}$ for $j<0$. The spherical harmonic space $\mathcal Y_k$ consists of restrictions of elements of $\mathcal H_k$ to the unit sphere $S^n$. Restriction is injective by homogeneity. On polynomial coefficients define the positive inner product
\[
 (p,q)_F=\sum_\alpha \alpha!p_\alpha q_\alpha,
 \qquad p(x)=\sum_\alpha p_\alpha x^\alpha.
\]
This coefficient inner product is not the surface $L^2$ inner product; it is an algebraic tool for the next decomposition.
\end{definition}

\begin{lemma}[Algebraic harmonic decomposition and dimensions]\label{lem:harmonic-polynomial-fischer-decomposition}
For $k\geq0$ one has
\[
 \mathcal P_k=\mathcal H_k\mathbin{\oplus}|x|^2\mathcal P_{k-2},
 \qquad
 \dim\mathcal Y_k=\binom{n+k}{k}-\binom{n+k-2}{k-2},
\]
where the second binomial term is zero for $k<2$. Iteration gives the unique expression $p=\sum_{j=0}^{\lfloor k/2\rfloor}|x|^{2j}h_{k-2j}$ with $h_{k-2j}\in\mathcal H_{k-2j}$. Compare \cite{ABR01}, Proposition 5.5, Theorem 5.7 and Proposition 5.8, pp.~76--78.
\end{lemma}
\begin{proof}
For monomials, multiplication by $x_i$ is adjoint to $\partial_i$ in the coefficient inner products: the factorial at index $\alpha+e_i$ is $(\alpha_i+1)\alpha!$. Thus multiplication by $|x|^2$ is adjoint to $\Delta$ between the indicated homogeneous spaces. Its range is therefore the orthogonal complement of $\ker\Delta=\mathcal H_k$ inside the finite-dimensional space $\mathcal P_k$. This proves the direct sum. Multiplication by $|x|^2$ is injective since the polynomial ring has no zero divisors, and $\dim\mathcal P_k=\binom{n+k}{k}$ by counting monomials. Subtraction gives the dimension formula. Repeatedly apply the unique direct-sum decomposition to the polynomial of degree $k-2$; the iteration stops at degree zero or one. This argument does not require a prior solution of a Dirichlet problem.
\end{proof}

\begin{proposition}[Eigenvalues and orthogonality of harmonic restrictions]\label{prop:harmonic-restriction-eigenvalues}
Every $Y\in\mathcal Y_k$ satisfies $-\Delta_{S^n}Y=\lambda_kY$, where $\lambda_k=k(k+n-1)$. Spaces with distinct degrees are orthogonal both in $L^2(S^n)$ and in the Dirichlet energy pairing.
\end{proposition}
\begin{proof}
In polar coordinates the Euclidean metric is $dr^2+r^2g_{S^n}$ and its volume element is $r^n\,dr\,d\sigma$. The divergence formula gives
\[
 \Delta_{\mathbb R^{n+1}}=\partial_r^2+\frac nr\partial_r+\frac1{r^2}\Delta_{S^n}.
\]
For a homogeneous polynomial $p(r\theta)=r^kY(\theta)$ this yields
$\Delta p=r^{k-2}(k(k+n-1)Y+\Delta_{S^n}Y)$. Harmonicity gives the equation. Integration by parts on the closed sphere then gives
$(\lambda_k-\lambda_j)\int Y_kY_j=0$ and $\int\langle\nabla Y_k,\nabla Y_j\rangle=\lambda_j\int Y_kY_j$. Since $\lambda_{k+1}-\lambda_k=2k+n>0$, both pairings vanish when $j\ne k$. This is the eigenfunction version of the orthogonality in \cite{ABR01}, Proposition 5.9, pp.~79--80.
\end{proof}

\begin{theorem}[Complete spherical harmonic expansion in L2]\label{thm:spherical-harmonic-l2-expansion}
Every real $f\in L^2(S^n)$ has a unique expansion $f=\sum_{k=0}^\infty f_k$ with $f_k\in\mathcal Y_k$, convergent in $L^2$, and $\|f\|_2^2=\sum_k\|f_k\|_2^2$. This is \cite{ABR01}, Theorem 5.12, p.~81, in real rather than complex notation.
\end{theorem}
\begin{proof}
Every polynomial is the finite sum of its homogeneous parts. Lemma \ref{lem:harmonic-polynomial-fischer-decomposition} shows that its restriction to the sphere lies in the finite span of the $\mathcal Y_k$, because $|x|=1$ there. Coordinate polynomials form a real algebra containing constants and separating points, so Stone--Weierstrass makes their restrictions uniformly dense in $C(S^n)$. Continuous functions are dense in $L^2$ for surface measure, and the finite total measure converts uniform approximation into $L^2$ approximation. Thus the harmonic span is $L^2$ dense. Each $\mathcal Y_k$ is finite dimensional and hence closed. Let $f_k$ be its orthogonal projection of $f$. Orthogonality shows that finite partial sums are the projections onto the corresponding finite spans. Density makes their residual norms tend to zero. Pythagoras gives the norm identity and uniqueness. No elliptic spectral existence theorem has been assumed.
\end{proof}

\begin{proposition}[Sobolev energy expansion and exhaustion of weak eigenfunctions]\label{prop:spherical-harmonic-h1-expansion}
If $f\in H^1(S^n)$, its preceding expansion converges in $H^1$ and
\[
 \int_{S^n}|\nabla f|^2=\sum_{k=0}^\infty\lambda_k\|f_k\|_2^2.
\]
Conversely an $L^2$ expansion with finite weighted sum on the right represents an $H^1$ function. If a nonzero $f\in H^1$ satisfies $\int\langle\nabla f,\nabla\phi\rangle=\lambda\int f\phi$ for every smooth $\phi$, then $\lambda=\lambda_k$ for one $k$ and $f\in\mathcal Y_k$.
\end{proposition}
\begin{proof}
For every smooth harmonic restriction $Y\in\mathcal Y_k$, weak integration by parts gives $\int\langle\nabla f,\nabla Y\rangle=\lambda_k\int fY$. If $S_K=\sum_{k\leq K}f_k$, the gradient of $f-S_K$ is therefore orthogonal to the gradient of $S_K$. Hence
\[
 \int|\nabla f|^2=\sum_{k\leq K}\lambda_k\|f_k\|_2^2+\int|\nabla(f-S_K)|^2.
\]
The nonnegative sums are bounded by the energy. Their tails, together with the $L^2$ norm tails, make $S_K$ Cauchy in $H^1$. By completeness it has an $H^1$ limit, whose $L^2$ value must be $f$ by Theorem \ref{thm:spherical-harmonic-l2-expansion}. This proves convergence and equality, avoiding the unsupported inference that $L^2$ convergence alone controls derivatives. The same Cauchy argument proves the converse. Finally test the weak eigenfunction equation against $Y\in\mathcal Y_k$: $(\lambda_k-\lambda)\int fY=0$. Completeness of the expansion implies that some projection is nonzero and that all other projections vanish. This identifies all weak eigenfunctions without asserting additional operator-domain facts.
\end{proof}

\begin{corollary}[Sharp spectral gaps after removing low modes]\label{cor:sphere-higher-mode-gap}
On $S^n(R)$, $R>0$, let $f\in H^1$ be orthogonal to the scaled spaces of degrees $0,\ldots,m-1$, where $m\geq1$. Then
\[
 \int|\nabla f|^2\geq\frac{m(m+n-1)}{R^2}\int f^2,
\]
with equality precisely for functions of degree $m$ (including zero). In particular removal of constants and linear modes gives the sharp factor $2(n+1)/R^2$.
\end{corollary}
\begin{proof}
On the unit sphere, Proposition \ref{prop:spherical-harmonic-h1-expansion} expresses the difference between the two sides as $\sum_{k\geq m}(\lambda_k-\lambda_m)\|f_k\|_2^2$. Strict increase of the eigenvalues makes equality equivalent to absence of every degree greater than $m$. Under $x=R\theta$, area scales by $R^n$ and the squared gradient by $R^{-2}$. This gives the factor $R^{-2}$ and preserves orthogonality. The assertion supplies the second positive gap used in Proposition \ref{prop:shrinker-sphere-f-stable}; it does not duplicate the already established first Poincare inequality.
\end{proof}

\begin{example}[Linear and trace-free quadratic modes]\label{ex:sphere-low-harmonic-modes}
The degree-zero space consists of constants. The degree-one space has dimension $n+1$ and consists of $x\mapsto a\cdot x$ restricted to the sphere. Degree two consists of $x\mapsto x^TBx$ with $B$ symmetric and $\operatorname{tr}B=0$; its dimension is $(n+1)(n+2)/2-1$. On the shrinker sphere $R=\sqrt{2n}$, the sharp inequality $Q(g)\geq [g^2]/n$ after removing constants and linear modes has equality exactly on this quadratic space. Here $[\,\cdot\,]$ denotes Gaussian integration on the sphere, and $Q$ is the existing Gaussian quadratic form.
\end{example}
\begin{proof}
All degree-one polynomials are harmonic. Every real homogeneous quadratic has a unique symmetric-matrix expression, and $\Delta(x^TBx)=2\operatorname{tr}B$, giving the stated characterization and dimension. On a sphere centered at the origin the Gaussian density is a positive constant, so weighted and ordinary orthogonality agree. At radius $\sqrt{2n}$, one has $|A|^2=1/2$ and therefore $Q(g)=[|\nabla g|^2-g^2]$. Corollary \ref{cor:sphere-higher-mode-gap} with $m=2$ gives $2(n+1)/(2n)-1=1/n$, with its equality characterization. On $S^1$, these degree-two modes are the familiar $\cos(2\theta)$ and $\sin(2\theta)$, so the dimension formula and equality case also include $n=1$.
\end{proof}

\begin{lemma}[Second and fourth coordinate moments]\label{lem:sphere-coordinate-moments}
Write $d=n+1\geq2$ and let $\langle f\rangle$ denote the surface average on the unit sphere in $\mathbb R^d$. Then
\[
 \langle x_i x_j\rangle=\frac{\delta_{ij}}d,\qquad
 \langle x_i x_j x_k x_l\rangle
 =\frac{\delta_{ij}\delta_{kl}+\delta_{ik}\delta_{jl}+\delta_{il}\delta_{jk}}{d(d+2)}.
\]
For symmetric trace-free matrices $B,C$ this gives
$\langle(x^TBx)(x^TCx)\rangle=2\operatorname{tr}(BC)/(d(d+2))$.
\end{lemma}
\begin{proof}
Reflection of any coordinate makes every monomial with an odd exponent in that coordinate integrate to zero. Permuting coordinates makes their second moments equal, and $\sum x_i^2=1$ gives the first formula. Put $a=\langle x_1^4\rangle$ and $b=\langle x_1^2x_2^2\rangle$. Rotational invariance makes $\langle((x_1+x_2)/\sqrt2)^4\rangle=a$, while expansion and odd-moment cancellation make it $a/2+3b/2$. Thus $a=3b$. Squaring $\sum x_i^2=1$ gives $da+d(d-1)b=1$, so $b=1/(d(d+2))$. The displayed tensor formula accounts for all possibilities of repeated indices. Contract it against $B_{ij}C_{kl}$: the three contractions are $\operatorname{tr}B\operatorname{tr}C$, $\operatorname{tr}(BC)$ and $\operatorname{tr}(BC)$, proving the final assertion.
\end{proof}

\begin{proposition}[Explicit projections onto the first three harmonic spaces]\label{prop:sphere-low-mode-projections}
For $f\in L^2(S^n)$ with the average notation above, its degree-zero, one and two projections are
\[
 f_0=\langle f\rangle,\qquad
 f_1(x)=d\langle f x\rangle\cdot x,\qquad
 f_2(x)=x^TB_f x,
\]
\[
 B_f=\frac{d(d+2)}2
       \left(\langle fxx^T\rangle-\frac{\langle f\rangle}{d}I\right).
\]
All averages exist by Cauchy--Schwarz on the finite measure sphere.
\end{proposition}
\begin{proof}
The formula for constants is their orthogonal projection. Lemma \ref{lem:sphere-coordinate-moments} gives $\langle x_i x_j\rangle=\delta_{ij}/d$, so the displayed linear function has the same pairing with each coordinate as $f$. The matrix $B_f$ is symmetric and trace free, since $\operatorname{tr}\langle fxx^T\rangle=\langle f\rangle$. For any symmetric trace-free $C$, the same lemma gives
\[
 \langle (x^TB_fx)(x^TCx)\rangle
 =\frac{2\operatorname{tr}(B_fC)}{d(d+2)}
 =\langle f(x^TCx)\rangle.
\]
Thus $f-f_2$ is orthogonal to the entire quadratic harmonic space and $f_2$ belongs to it, proving the formula. These are actual orthogonal projections, not coefficients for an assumed finite polynomial approximation to $f$.
\end{proof}

\begin{proposition}[A mean-zero weak Poisson inverse on a round sphere]\label{prop:sphere-poisson-inverse}
For $f\in L^2(S^n(R))$ with zero mean there is a unique mean-zero $u\in H^1$ such that
\[
 \int\langle\nabla u,\nabla\phi\rangle=\int f\phi
 \quad\text{for every }\phi\in H^1.
\]
Writing $\lambda_k=k(k+n-1)/R^2$, it is given by the $H^1$ convergent series $u=\sum_{k\geq1}f_k/\lambda_k$, and
\[
 \|u\|_2\leq\frac{R^2}{n}\|f\|_2,\qquad
 \|\nabla u\|_2\leq\frac R{\sqrt n}\|f\|_2.
\]
This is a weak inverse statement; higher regularity is not being inferred solely from $L^2$ convergence of a series.
\end{proposition}
\begin{proof}
The eigenvalues for $k\geq1$ are at least $\lambda_1=n/R^2$. Hence
\[
 \sum_{k\geq1}\frac{\|f_k\|_2^2}{\lambda_k^2}
 \leq\lambda_1^{-2}\|f\|_2^2,\qquad
 \sum_{k\geq1}\frac{\|f_k\|_2^2}{\lambda_k}
 \leq\lambda_1^{-1}\|f\|_2^2.
\]
Proposition \ref{prop:spherical-harmonic-h1-expansion}, with the scaling in Corollary \ref{cor:sphere-higher-mode-gap}, proves $H^1$ convergence and the two bounds. The weak identity holds for finite harmonic sums by their eigenfunction equations. Those sums converge in $H^1$ for any test function, so Cauchy--Schwarz passes both sides to the limit. If two mean-zero solutions existed, test the equation for their difference against that difference. Its gradient has zero norm; the positive first gap makes the function zero. This gives uniqueness without a separate elliptic existence theorem.
\end{proof}
